\section{模型总结}

本文围绕微网在不同信息条件下的购电与储能调度问题，构建了由日前计划、预测场景与日内修正逐层递进的统一框架。问题一在电价、负荷和光伏完全已知时，以电力平衡、储能状态递推和充放电互斥为约束，给出确定性优化基准；问题二将信息边界限定在每日 0:00，以周期加权预测和整日联合残差场景刻画负荷、光伏的不确定性，并以季节余弦终端软区间维持跨日储能价值；问题三利用官方光伏预报在 6:00、12:00、18:00 对剩余时域滚动优化；问题四进一步处理实时波动电价，以周期预测为基线、以因果门控残差模型完成日内更新。四个问题的约束内核一致，因而不同信息结构下的结果具有可比性。

结果表明，储能能够通过低价充电、高价放电和吸收光伏富余电量有效降低购电成本：完全信息的确定性调度相对无储能基准节省 26.90\%，并消除弃光。面对负荷、光伏及电价的不确定性，严格按可得信息制定的日前与滚动策略仍保持稳定收益。问题三的四时点策略较仅 0:00 决策节省 2.36\%，问题四 4-3 的对应降幅为 2.44\%，且两类情形均显示 6:00、12:00 更新最具经济价值。价格完全信息下界与现实策略的费用差距约为 1\%，说明周期预测结合门控更新已提取了大部分可用于储能套利的价格信息。

本模型的优点在于：预测、残差归档、场景抽样和更新决策均遵循严格的时间信息边界；储能容量、功率、效率及充放电互斥约束保证策略可执行；以实际数据回放结算，使费用结论可复核、可复现。模型尚未显式计入电池退化成本、向外部电网售电机制及极端天气导致的预测分布突变，且终端软区间仍为经验性跨日备用安排。后续可引入寿命损耗与碳排放等多目标成本，结合更长预测窗口和鲁棒或分布鲁棒优化，进一步提升极端场景下的经济性与安全裕度。
